%% Chapter 17 ----------------------------------------------------------------
\chapter{Image Display and Processing}
\label{ch:solar-processing}

This chapter describes the tools of Solar Image Analysis that display, enhance, combine and
export image sequences once they are loaded (Chapter~\ref{ch:solar-window}).

\section{The Analysis card}
\label{sec:solar-analysis-card}

The \gui{Analysis} card (Figure~\ref{fig:solar-cards}a) shows a summary of the loaded data and
the main display modes and analysis tools.

\begin{figure}[!htbp]
  \centering
  \begin{minipage}[b]{0.47\linewidth}\centering
    \includegraphics[width=\linewidth]{figures/solar/solar_analysis_cards.png}\par\smallskip
    {\small (a) \gui{Analysis} card}
  \end{minipage}\hfill
  \begin{minipage}[b]{0.47\linewidth}\centering
    \includegraphics[width=\linewidth]{figures/solar/solar_display_cards.png}\par\smallskip
    {\small (b) \gui{Display \& Crop} card}
  \end{minipage}
  \caption{The \gui{Analysis} and \gui{Display \& Crop} cards of Solar Image Analysis.}
  \label{fig:solar-cards}
\end{figure}

\begin{description}
  \item[Plot Frames] displays the loaded sequence as plain frames.
  \item[Running Diff] displays each frame minus the previous one, which makes moving and
    changing features stand out.
  \item[Base Diff] displays each frame minus the first, showing the total change since the
    start --- useful for eruptions, coronal dimmings and EUV waves.
  \item[Composite] builds an RGB tri-color image from the first three AIA frames or, when a
    magnetogram overlay has been loaded, draws HMI polarity contours on the current AIA frame
    (red positive, blue negative).
  \item[Mag Overlay…] loads an HMI magnetogram FITS file for the composite; the box beside it
    sets the contour level (10--2000~G; default 100~G).
  \item[Light Curve] plots the mean intensity of a region against time
    (Section~\ref{sec:region-lightcurve}).
  \item[Active Regions], \gui{NOAA/HEK Labels}\index{active regions!NOAA labels} detect and label active regions
    (Section~\ref{sec:active-regions}).
  \item[Compare Viewpoint…] compares the frame with a second observer
    (Section~\ref{sec:compare-viewpoint}).
  \item[Reset Frames] restores the frames as downloaded, discarding crops, processing
    levels, derotation, composites and detected regions.
\end{description}
\gui{Composite}\index{composite image}, \gui{Mag Overlay…}\index{magnetogram overlay}, \gui{Active Regions} and \gui{NOAA/HEK Labels} are
available only for disk imagers and magnetographs. The difference modes also determine the
content of exported movies (Section~\ref{sec:solar-movies}).

\section{Display and cropping}
\label{sec:display-crop}
\index{display!solar images}\index{crop}

The \gui{Display \& Crop} card (Figure~\ref{fig:solar-cards}b) controls the appearance of the image:
\begin{description}
  \item[Renderer] \gui{PyQtGraph}\index{PyQtGraph} (default; fast, and required for clicked PFSS seeds and the
    interactive crop rectangle) or \gui{Matplotlib}\index{Matplotlib}.
  \item[Colormap] instrument color maps (for example \texttt{sdoaia193}, \texttt{hmimag},
    \texttt{soholasco2}, \texttt{stereocor2}, \texttt{euvi195}, \texttt{goes-rsuvi171}) and
    general ones (\texttt{inferno}, \texttt{magma}, \texttt{plasma}, \texttt{viridis},
    \texttt{cividis}, \texttt{gray}, \texttt{RdBu\_r}, \texttt{hot}). A suitable map is
    chosen automatically for each observable.
  \item[Scale] \gui{linear} or \gui{log}.
  \item[Clip low, Clip high] the lower and upper percentiles mapped to the ends of the color
    scale (defaults 1\% and 99.9\%); dragging gives a live preview.
  \item[Solar Limb] draws the photospheric limb.
  \item[Coordinate Grid, Grid frame] overlays a graticule of meridians and parallels in the
    selected frame --- Heliocentric Inertial (HCI, the default), Stonyhurst or Carrington. The
    same frame is used for the heliographic coordinates of the cursor readout.
  \item[Colorbar, Region Overlays] show the color bar and the detected active regions.
\end{description}

\procedure{Cropping}
Under \gui{Crop (arcsec)}, enter \gui{X min / max} and \gui{Y min / max} in arcseconds from
disk center, or tick \gui{Rectangle crop}\index{crop!rectangle} to drag a rectangle on the image that fills in
the fields. \gui{Apply Crop} crops all loaded frames to the rectangle. Cropping is done
locally after the files are loaded. \gui{Export Cropped FITS}\index{crop!export FITS} saves the cropped frames as
FITS files with updated coordinate metadata.

The crop rectangle also defines the region used by \gui{Region Stats}\index{measurement tools!region statistics}, the light curve, the
derotation and the \gui{Current crop} PFSS seeding.

\section{Solar derotation}
\label{sec:derotation}
\index{derotation}\index{differential rotation}

Over a sequence of several hours, a fixed rectangle slowly drifts onto neighboring
surface because of solar differential rotation. The \gui{Solar Derotation} card keeps a
selected region on the same piece of the Sun:
\begin{description}
  \item[Derotate selected area] enables derotation of the crop region.
  \item[Mode] \gui{Track region (no resampling)}\index{derotation!track mode} moves the cut-out window with the rotation
    and keeps the original pixel values, so it is the mode for photometry and light curves.
    \gui{Reproject to reference time}\index{derotation!reproject mode} resamples every frame onto the reference time's grid,
    which corrects foreshortening and gives all frames one coordinate system, making
    differencing meaningful, at the cost of interpolating the data.
  \item[Reference] \gui{First frame}, \gui{Current frame} or \gui{Specific frame...} (with its
    number): the frame from which the region is read and to whose time every other frame is
    rotated.
  \item[Padding] the extra margin cut around the region before reprojecting (0--600~arcsec;
    default 60~arcsec), so that the interpolation has data up to the edge. Reproject mode only.
  \item[Apply Derotation, Reset to Full Frames] apply the derotation, or discard the cut-out
    and restore the full frames at the current processing level.
\end{description}
The status line reports whether derotation is active.

\section{Region light curve}
\label{sec:region-lightcurve}
\index{light curve!EUV}

\gui{Light Curve} plots the mean intensity of the crop region (or of the whole frame if
\gui{Rectangle crop} is off), converted to DN/s using the exposure time, against time
(Figure~\ref{fig:region-lightcurve}). The time of peak intensity is marked and, when a
spectrum is loaded in the main window, its time window is shaded so that the EUV brightening
can be timed against the radio burst. At least two frames are needed.

\screenshot[0.72\linewidth]{solar/region_light_curve}{The \gui{Region Light Curve} window.}{fig:region-lightcurve}{The Region Light Curve window for an AIA flare region, showing the DN/s curve, the EUV peak marker and the shaded radio burst window.}

\section{Active regions}
\label{sec:active-regions}
\index{active regions}

The \gui{Active Regions} card detects bright regions in the current frame and lists them:
\begin{description}
  \item[Threshold] the intensity percentile above which pixels count as bright (50--100;
    default 98).
  \item[Min px] the minimum area of a region in pixels (default 12).
  \item[Region table] one row per region with \gui{ID}, \gui{Label}, \gui{NOAA} number,
    \gui{Centroid}, \gui{Area}, \gui{Peak} intensity and \gui{Source}.
\end{description}
Click \gui{Active Regions} in the \gui{Analysis} card to run the detection, and
\gui{NOAA/HEK Labels} to fetch the NOAA active-region numbers from the Heliophysics Event
Knowledge base (HEK) and attach them to the detected regions; the status line reports whether
metadata are loaded. Detected regions are drawn on the image when \gui{Region Overlays}\index{active regions!overlays} is
ticked. \gui{Export Regions CSV} saves the table with centroids, bounding boxes and
intensities. Detection works on local files; the labels need network access.

\section{Comparing viewpoints}
\label{sec:compare-viewpoint}
\index{Compare Viewpoint}

\gui{Compare Viewpoint…} opens a window that compares the current frame (viewpoint~A) with a
second observable (viewpoint~B) taken near the same time --- for example a STEREO image
against the Earth view.
\begin{steps}
  \item Choose the \gui{Second viewpoint} observable and the search window (\gui{±}, 1--720~min;
    default 30~min).
  \item Click \gui{Fetch Viewpoint B}. Viewpoint~B is downloaded, reprojected onto the
    coordinate frame of A, and shown beside it.
  \item Tick \gui{Blink}\index{Compare Viewpoint!blink} to alternate between the two images in place.
\end{steps}
The \gui{Observer separation}\index{Compare Viewpoint!observer separation} line gives the angle between the two observers.

\section{Coronagraph tools}
\label{sec:nrgf}
\index{NRGF}

For coronagraph data the \gui{Coronagraph Tools} card offers \gui{NRGF radial filter}\index{NRGF}, the
normalizing-radial-graded filter \parencite{morgan2006}. It flattens the steep radial fall-off
of coronal brightness so that faint CME fronts become visible at all heights. It applies to
the plain frame view; difference images already remove the background.

\section{Overlay layers}
\label{sec:overlay-layers}
\index{overlay layers}\index{coronagraph composite}

A CME is never fully visible in one instrument: the eruption starts on the disk (AIA, EIT,
EUVI), crosses the low corona (LASCO~C2, COR1) and expands into the outer corona (LASCO~C3,
COR2), and each coronagraph is blind inside its occulter. For coronagraph views the
\gui{Overlay Layers} card blends several instruments into one continuous image from the disk
to the outer corona.

\begin{description}
  \item[Add, Remove] choose an observable and add it to the layer list, or remove the
    selected layer. Layers are painted widest field of view first. Untick a layer to leave it
    out without losing its settings.
  \item[Per-layer settings] for the selected layer: \gui{Colormap}, \gui{Scale} (log or
    linear), \gui{Opacity}, \gui{Midtones}\index{overlay layers!midtones} (gamma; 50 suits most solar images), and the
    \gui{Inner} and \gui{Outer} edges of its field of view in \(R_\odot\). Pixels inside the
    inner edge stay transparent so that the layer below shows through.
  \item[Time match] how far from each loaded frame's time a layer frame may be taken
    (1--720~min; default 30~min). Widen it for slow cadences or patchy archives.
  \item[Build Composite] downloads each layer's matching frames, reprojects them onto the
    loaded series and blends them, in the background with a progress bar. The result replaces
    the view and works with the measurement tools, plot export and movie export.
  \item[Clear] drops the composite and restores the originally loaded frames.
\end{description}

\begin{note}
  Each layer is color-mapped before blending, because EUV disk images and the white-light
  corona differ by orders of magnitude. Coronagraph layers are reprojected under a
  spherical-screen assumption, flagged as approximate in the build notes: cross-observer
  coronagraph overlays are morphological context, not photometry. SOHO/LASCO\index{SOHO/LASCO} files carry no
  observer position, so an Earth-based observer is assumed (accurate to about 1\%).
\end{note}

\section{Heliospheric imager J-maps}
\label{sec:jmap}
\index{J-map}\index{heliospheric imager}

For STEREO HI1 and HI2 data the \gui{Heliospheric Imager (J-map)} card builds a
time--elongation map \parencite{sheeley1999}:
\begin{description}
  \item[Background] \gui{Median background} removes the static F-corona and star field with
    the per-pixel temporal median; \gui{Previous frame} uses a running difference\index{running difference}.
  \item[PA] the position angle of the slit, counter-clockwise from the image \(+x\) axis
    (90° points up; default 90°). Set it along the CME direction.
  \item[Build J-map] stacks the slit profiles of all frames into the \gui{HI J-map} window
    (elongation in degrees against frame index). An outward-moving CME appears as a slanted
    bright track.
\end{description}

\section{HMI vector magnetic field}
\label{sec:hmi-vector}
\index{HMI!vector field}

For SDO/HMI\index{SDO/HMI} data the \gui{Magnetic Vector Field (HMI)} card overlays the measured
photospheric vector field:
\begin{description}
  \item[Load Vector FITS…] loads the \texttt{hmi.B\_720s} field, inclination and azimuth
    segments of each time step from disk (with the optional disambiguation segment, applied
    automatically).
  \item[Download Vector (JSOC)] downloads \texttt{hmi.B\_720s} for the query time window
    through JSOC\index{JSOC}; this needs the registered JSOC e-mail. Each 720-s time step is four segments
    of about 50~MB. If no HMI image is loaded, the derived vertical-field (\(B_z\))
    magnetogram is plotted with the vectors.
  \item[Show vector field] overlays the field on the HMI frame. The transverse component is
    drawn as \gui{Arrows} (length proportional to strength) and/or \gui{Streamlines}\index{HMI!streamlines}, and
    \gui{|B| layer} tints the image by total field strength. Arrows are red where the vertical
    field points toward the observer (\(+B_z\)) and blue where it points away.
  \item[Spacing, Threshold] the arrow grid spacing (16--512 detector pixels; default 64) and
    the minimum transverse field drawn (50--2000~G; default 200~G). The transverse noise floor
    of \texttt{hmi.B\_720s} is about 100~G.
\end{description}

\section{Movies and still images}
\label{sec:solar-movies}
\index{movie export}

\begin{description}
  \item[Movie Export card] choose the \gui{Content} (\gui{Frames}, \gui{Running Difference}
    or \gui{Base Difference}\index{base difference}) and the \gui{Format} (\gui{MP4} or \gui{GIF}), and click
    \gui{Build Movie…}\index{movie export}. Movies use the playback speed set under the viewer and the current
    color map, scale, clipping and crop.
  \item[Export MP4] in the header bar makes a one-click MP4 of the loaded sequence with the
    current display settings.
  \item[Export Plot] saves the current frame, with its overlays, as an image.
  \item[Export Cropped FITS] saves the cropped frames as FITS files.
  \item[Export Regions CSV] saves the active-region table.
\end{description}
The same exports are available from the \menu{Export} and \menu{Movie} menus.
